%% -------------------------------------------------------
%% Kapitel 4: Technologieauswahl
%% -------------------------------------------------------
\chapter{Technologieauswahl}
\label{chap:technologie}

Die Wahl geeigneter Technologien hat einen entscheidenden Einfluss auf Wartbarkeit,
Performance und Entwicklungsgeschwindigkeit eines Projekts. Dieses Kapitel begründet die
getroffenen Entscheidungen für Backend und iOS-App und bewertet jeweils relevante
Alternativen.

\section{Backend-Framework: ASP.NET Core 8}
\label{sec:aspnet}

\subsection{Bewertete Alternativen}
\label{subsec:backend-alternativen}

Für die Implementierung der REST-API wurden drei Backend-Frameworks in Betracht gezogen:

\begin{table}[H]
  \centering
  \caption{Vergleich Backend-Frameworks}
  \label{tab:backend-vergleich}
  \small
  \begin{tabularx}{\textwidth}{|l|Y|Y|Y|}
    \hline
    \textbf{Kriterium} & \textbf{ASP.NET Core 8} & \textbf{Node.js / Express} &
    \textbf{Spring Boot 3} \\
    \hline
    Sprache             & C\# (.NET)      & JavaScript / TypeScript & Java / Kotlin \\
    \hline
    Performance         & Sehr hoch (Benchmark: Top-3 TechEmpower) &
                          Hoch (Event Loop) & Hoch \\
    \hline
    ORM-Unterstützung   & Entity Framework Core (erstklassig) &
                          Prisma / TypeORM & Hibernate / Spring Data \\
    \hline
    Lernkurve           & Mittel (C\# Vorwissen vorhanden) &
                          Niedrig (JS bekannt) & Hoch \\
    \hline
    Docker-Support      & Offizielles .NET-Image, klein & Node-Image & JVM-Image, groß \\
    \hline
    Schul-Kontext       & HTL-Lehrplan (C\#) & Partiell & Nein \\
    \hline
  \end{tabularx}
\end{table}

\subsection{Entscheidung}
\label{subsec:aspnet-entscheidung}

\textbf{ASP.NET~Core~8} wurde gewählt, weil C\# im HTL-Lehrplan fest verankert ist und
bereits umfangreiche Vorkenntnisse vorhanden sind. Das Framework bietet erstklassige
Integration mit Entity~Framework~Core, sehr gute Performance und einen klar strukturierten
Middleware-Stack. Die offiziellen Microsoft-Docker-Images sind schlank und werden langfristig
gewartet. Spring~Boot wäre eine technisch gleichwertige Alternative, erfordert jedoch Java
und eine deutlich steilere Einarbeitungskurve in das Spring-Ökosystem. Node.js~/ Express
wurde verworfen, da der Typsicherheitsgewinn gegenüber JavaScript und die bessere
ORM-Integration von EF~Core für ein datenbankzentrisches Projekt ausschlaggebend war.

\section{ORM: Entity Framework Core 9}
\label{sec:efcore}

\subsection{Bewertete Alternativen}
\label{subsec:orm-alternativen}

Für den Datenbankzugriff wurden drei Ansätze verglichen:

\begin{description}
  \item[Entity Framework Core 9 (EF Core)] Vollständiges ORM mit Code-First-Migrations,
    LINQ-Abfrageunterstützung und Change~Tracking.
  \item[Dapper] Lightweight Micro-ORM: Nur SQL-Mapping, kein Change~Tracking, keine
    automatischen Migrations.
  \item[ADO.NET (direkt)] Rohe Datenbankzugriffe; maximale Kontrolle, aber hoher
    Boilerplate-Aufwand.
\end{description}

\subsection{Entscheidung}
\label{subsec:efcore-entscheidung}

EF~Core~9 wurde bevorzugt, weil es automatische Migrationen ermöglicht (kein manuelles
SQL für Schemaänderungen), LINQ-Abfragen typsicher und intellisense-fähig macht und
sehr gut mit ASP.NET~Core zusammenspielt. Für einfache CRUD-Operationen wie in DoggoGO
ist der Overhead von EF~Core vernachlässigbar. Dapper wäre bei komplexen, performance-kritischen
Datenbankabfragen sinnvoller; da DoggoGO aktuell keine komplexen Joins oder Massenoperationen
benötigt, ist dieser Vorteil irrelevant.

\section{Datenbank: PostgreSQL}
\label{sec:postgresql}

\subsection{Bewertete Alternativen}
\label{subsec:db-alternativen}

\begin{table}[H]
  \centering
  \caption{Vergleich Datenbanksysteme}
  \label{tab:db-vergleich}
  \small
  \begin{tabularx}{\textwidth}{|l|Y|Y|Y|}
    \hline
    \textbf{Kriterium} & \textbf{PostgreSQL} & \textbf{MySQL / MariaDB} & \textbf{SQLite} \\
    \hline
    ACID-Konformität  & Vollständig & Vollständig & Vollständig \\
    \hline
    Array-Typ (text[]) & Nativ unterstützt & Nicht nativ & Nicht vorhanden \\
    \hline
    Docker-Image      & Offiziell, stabil & Offiziell & Nicht nötig \\
    \hline
    Skalierbarkeit    & Hoch & Hoch & Niedrig (Datei-basiert) \\
    \hline
    EF Core Support   & Npgsql (erstklassig) & Pomelo (gut) & Eingebaut \\
    \hline
  \end{tabularx}
\end{table}

\subsection{Entscheidung}
\label{subsec:postgresql-entscheidung}

\textbf{PostgreSQL} wurde gewählt, weil das Modell \texttt{LegalRestrictions} die Regeln
als \texttt{List<string>} speichert – ein Feature, das in PostgreSQL nativ als
\texttt{text[]} (Array-Typ) unterstützt wird und von Npgsql direkt auf C\#-Listen
gemappt werden kann. MySQL~/ MariaDB bietet diesen nativen Array-Typ nicht; dort wäre
eine separate Normalisierungstabelle oder JSON-Serialisierung notwendig gewesen, was den
Code komplexer gemacht hätte. SQLite schied aus, da die Datei-basierte Architektur für
eine containerisierte Multi-User-Anwendung nicht geeignet ist.

\section{iOS: Swift und SwiftUI}
\label{sec:swiftui}

\subsection{Bewertete Alternativen}
\label{subsec:ios-alternativen}

Für die native iOS-App wurden drei Ansätze verglichen:

\begin{description}
  \item[Swift / SwiftUI (nativ)] Offizieller Apple-Stack, maximale Plattformintegration,
    beste MapKit- und Keychain-Unterstützung.
  \item[React Native] JavaScript-basiert, Cross-Platform (iOS + Android). MapKit-Unterstützung
    über Drittbibliotheken, weniger tief integriert.
  \item[Flutter] Dart-basiert, Cross-Platform. Google-Maps-Integration gut; Apple-Ökosystem
    (Keychain, HealthKit etc.) schwieriger anzubinden.
\end{description}

\subsection{Entscheidung}
\label{subsec:swiftui-entscheidung}

Da das Projektziel explizit eine \textbf{iOS}-Applikation vorsieht (kein Android-Support
im Projektrahmen) und MapKit sowie der iOS~\texttt{Keychain} für Kartendarstellung
und sichere Token-Speicherung zentrale Anforderungen sind, fiel die Wahl auf den
nativen \textbf{Swift~/ SwiftUI}-Stack. SwiftUI ermöglicht eine deklarative UI-Entwicklung
mit weniger Boilerplate als UIKit und integriert sich nahtlos in das Apple-Ökosystem.
Cross-Platform-Frameworks wie React~Native oder Flutter bieten für reine iOS-Anforderungen
keinen Mehrwert, führen aber zu einer tieferen Abstraktionsschicht und eingeschränkterem
Zugang zu nativen APIs.
